The anchor for this work is the Alignment Forum post ``Data filtering
works a lot worse than you would expect'' \citep{lee2026filtering}. It
found that when you fine-tune a base model on chat data, some of the
behaviors this installs can be removed by data filtering and some cannot.
Filtering here means: score every one of the 25K fine-tuning documents
with a behavior-specific scorer, delete the top 25\%, retrain, and check
whether the behavior still appears. Refusal mostly disappears under this
treatment; both-sides framing survives it completely. A companion informal
research update narrows the causal question by separating prompt
distribution, visible completion content, and teacher model or rollout source
\citep{engels2026filters}; it does not settle a general mechanism. The
question matters for any safety strategy that leans on training-data
curation: if we cannot predict which behaviors filtering can touch, we
cannot rely on it. I test one proposed persona account of this split rather
than treating the filtering contrast as an already-solved causal result.
